% 解説記事：特集内の記事扉．左寄せの見出し，副題，飾り，段抜きの囲みを使う形．
% 本文は本例のために書き下ろした説明文であり，実在の記事の引用ではない．
\documentclass[lualatex,paper=b5j,fontsize=10pt,twoside]{jlreq}

\usepackage{surikumi-kaisetsu}

% 記事扉の飾り．参考紙面の図案を写さず，TeX で生成した独自の図形を置く．
\newcommand{\MyOrnament}{%
  \begin{tikzpicture}[x=1mm,y=1mm,line width=.45pt]
    \fill[SMink] (0,0) circle (.9);
    \fill[SMink] (30,10) circle (.9);
    \draw[SMrule] (0,0) .. controls (10,14) and (20,14) .. (30,10);
    \draw[SMrule] (0,0) .. controls (12,6) and (18,2) .. (30,10);
    \draw[SMrule] (0,0) .. controls (8,-8) and (24,-4) .. (30,10);
    \draw[SMink,line width=.8pt] (0,0) .. controls (12,1) and (20,7) .. (30,10);
  \end{tikzpicture}%
}

\MathKaisetsuSetup{
  journal-name={数理組版},
  issue-date={February 2026},
  issue-number={Number 2},
  % 値に読点を含む鍵は，全体をもう一重の中括弧で囲む．
  footer-line={{\MKHeadingFace 数理組版}\hspace{.7em}NO. 2, FEBRUARY 2026},
  align=left,
  masthead=false,
  feature-bar=true,
  section-style=pillars,
  feature={特集／積分の考え方},
  title={面積を測るという操作},
  subtitle={Riemann 積分から測度へ},
  author={\MKName{組版}{花子}},
  author-reading={くみはん・はなこ},
  affiliation={数理組版編集部},
  ornament={\MyOrnament},
  ornament-width={34mm}
}

\begin{document}

\setcounter{page}{12}

\MKMakeTitle
\MKArticleReset

面積を求める操作は，長方形の面積からはじめて，
細かく分けて足し合わせるという手続きで拡張されてきた．
この手続きをどこまで押し進められるかを問うと，
積分の定義そのものを選び直す必要が生じる．
本稿では，分割の仕方を変えるという一つの発想が
Riemann 積分から Lebesgue 積分への移行を導く様子をたどる．

\MKSection{分割と近似}

区間を細かく分ける，各断片で高さを一つ選ぶ，
断片の幅と高さの積を足す．
Riemann 積分の定義はこの三つの操作からできている．

\MKSubsection{Riemann 和}

閉区間 $[a,b]$ の分割
$a=x_0<x_1<\cdots<x_n=b$ と，
各小区間の代表点 $\xi_k\in[x_{k-1},x_k]$ に対して
\begin{equation}
  R(f;\Delta,\xi)
    =\sum_{k=1}^{n}f(\xi_k)\,(x_k-x_{k-1})
\end{equation}
を Riemann 和という．
分割の幅の最大値を $\lVert\Delta\rVert$ と書き，
$\lVert\Delta\rVert\to0$ のとき $R(f;\Delta,\xi)$ が
代表点の取り方によらず一つの値へ近づくならば，
$f$ は $[a,b]$ で積分可能であるという\MKCite{1}．

\begin{figure}[t]
  \begin{mkfigurebody}
    \begin{tikzpicture}[x=1mm,y=1mm,line width=.5pt]
      \draw[SMrule,line width=.4pt] (0,0) -- (52,0);
      \draw[SMrule,line width=.4pt] (0,0) -- (0,28);
      \foreach \i in {0,...,5}{%
        \pgfmathsetmacro{\xa}{8*\i}%
        \pgfmathsetmacro{\ya}{6+16*sin((\xa+6)*1.7)*sin((\xa+6)*1.7)}%
        \draw[SMrule,line width=.35pt,fill=SMpale]
          (\xa,0) rectangle (\xa+8,\ya);%
      }
      \draw[SMink,line width=.7pt,domain=0:48,samples=80,smooth]
        plot (\x,{6+16*sin((\x+2)*1.7)*sin((\x+2)*1.7)});
      \node[font=\fontsize{6.5}{7}\selectfont] at (-2.6,26) {$y$};
      \node[font=\fontsize{6.5}{7}\selectfont] at (51,-2.6) {$x$};
      \node[font=\fontsize{6.5}{7}\selectfont] at (0,-2.6) {$a$};
      \node[font=\fontsize{6.5}{7}\selectfont] at (48,-2.6) {$b$};
    \end{tikzpicture}
  \end{mkfigurebody}
  \caption{分割の各小区間で高さを一つ選び，長方形の面積を足し合わせる．
    分割を細かくしたときの極限が Riemann 積分である．}
\end{figure}

\MKSubsection{上限和と下限和}

代表点の選び方をすべて試す代わりに，
各小区間での上限と下限を取れば
\begin{equation}
  s(f;\Delta)\leq R(f;\Delta,\xi)\leq S(f;\Delta)
\end{equation}
がつねに成り立つ．
分割を細かくすると $s$ は増加し $S$ は減少するから，
両者が同じ値に収束することが積分可能性の判定条件になる．
連続関数がつねに積分可能であることは，
一様連続性からこの判定条件を確かめて示される\footnote{閉区間上の
連続関数は一様連続である．この事実を使うと，
分割を十分細かく取れば $S-s$ をいくらでも小さくできる．}．

\MKSection{分割の向きを変える}

Riemann 和では定義域を分割した．
値域を分割すると，同じ面積が別の形で書ける．

\MKSubsection{値域による分割}

$f$ の値域を $0=y_0<y_1<\cdots<y_m$ と分け，
$E_j=\{x\mid y_{j-1}\leq f(x)<y_j\}$ と置く．
このとき和
\begin{equation}
  \sum_{j=1}^{m}y_{j-1}\,\mu(E_j)
\end{equation}
を考える．ここで $\mu(E_j)$ は集合 $E_j$ の大きさであり，
区間であれば長さそのものである．
値域の分割を細かくした極限が Lebesgue 積分である\MKCite{2}．

\begin{mkremark}
定義域の分割では，各小区間で関数が大きく振れると
上限和と下限和が離れてしまう．
値域の分割では，関数の振れ方ではなく
集合 $E_j$ の測りやすさが問題になる．
\end{mkremark}

\MKSubsection{測度という道具}

区間の長さを出発点にして，
可算個の和と補集合について整合する形で
集合の大きさを定めたものが測度である．
測度が定まる集合の族を可測集合族という．

\begin{mkquote}
長さを測れる集合とは何か．
この問いに答えることが，積分を定義することにあたる．
\end{mkquote}

すべての集合に長さを与えることはできない．
可測でない集合が存在するという事実は，
測度をどこまで広げられるかの限界を示している\MKCite{3}．

\MKSection[band]{二つの積分の関係}

Riemann 積分可能な関数は Lebesgue 積分可能であり，
値も一致する．逆は成り立たない．
有理数の点で $1$，無理数の点で $0$ をとる関数は，
Riemann 和が代表点の取り方で変わるため Riemann 積分できないが，
Lebesgue 積分では値が $0$ と定まる．

\begin{mkinsert}
\MKRunIn{要点}
定義域を分ける立場では，関数の連続性が効いてくる．
値域を分ける立場では，集合の可測性が効いてくる．
同じ面積を測るという操作でも，
どこを細かく分けるかによって必要な道具が変わる．
\end{mkinsert}

極限と積分の交換が自由になることが，
値域を分割する立場の実際的な利点である．
単調収束定理と優収束定理は，
どちらも Riemann 積分では扱いにくかった状況を扱えるようにする．

\MKSubsection{関数列の極限}

各点で $f_n(x)\to f(x)$ となる関数列を考える．
Riemann 積分の枠内では，極限の関数が積分可能である保証がない．
値域を分割する立場では，$f_n$ が単調に増加していれば
\begin{equation}
  \lim_{n\to\infty}\int f_n\,\mathrm{d}\mu=\int f\,\mathrm{d}\mu
\end{equation}
がそのまま成り立つ．
積分の順序と極限の順序を入れ替えてよい条件が，
簡潔な形にまとまる点が要になる．

\MKSubsection{測ることの限界}

長さの概念を平行移動で不変な形に保ちながら
すべての部分集合へ広げることはできない．
この事実は，測度を定義するときに
可測集合の族をあらかじめ限っておく理由になっている．
積分を広げる操作には，
どこまで広げるかという選択が必ず伴う．

\begin{mkfullinsert}[b]
\MKRunIn{この特集について}
積分をめぐる三つの見方――区分求積，測度，超関数――を
それぞれ独立に読める記事として並べた．
どの記事も高校数学の範囲から書き起こしてあるので，
関心のある順に読み進めてよい．
\end{mkfullinsert}

\begin{mkreferences}
  \item 区分求積から積分へ進む標準的な微分積分学の教科書を参照した．
  \item 測度論の入門書は，区間の長さの拡張として測度を導入している．
  \item 可測でない集合の構成には選択公理を用いる．
\end{mkreferences}

\MKByline

\end{document}
